% This is text associated with the open textbook "Open Electromagnetics"
% See immediately below for author, date
% This work is licensed under the Creative Commons Attribution-ShareAlike 4.0 International License. (CC BY SA 4.0) 
% To view a copy of the license, visit http://creativecommons.org/licenses/by-sa/4.0/

%ORB TITLE Conductivity of Some Common Materials
%ORB AUTHOR 0        # S.W. Ellingson, Virginia Tech (ellingson@vt.edu)
%ORB DATE 20191128   # yyyymmdd
%ORB ABSTRACT Conductivity of Some Commonly-Encountered Materials Materials

\label{m0137_Table_of_Conductivities} % <-- Must match name of this file and module directory name.  Don't mess with this!
\index{conductivity!of common materials|(}

The values below are conductivity $\sigma$ for a few materials that are commonly encountered in electrical engineering applications, and for which conductivity emerges as a consideration.

Note that materials in some applications are described instead in terms of \emph{resistivity}, 
\index{resistivity}
which is simply the reciprocal of conductivity.

Conductivity may vary significantly as a function of frequency.
The values below are representative of frequencies from a few kHz to a few GHz.
Conductivity also varies as a function of temperature.
In applications where precise values are required, primary references accounting for frequency and temperature should be consulted.
The values presented here are gathered from a variety of references, including those indicated in ``Additional References'' at the end of this section.
  
{\bf Free Space} (vacuum): $\sigma\triangleq 0$.

{\bf Commonly encountered elements}:\\
%
\begin{tabular}{ll}
Material                   & $\sigma$ (S/m) \\
\hline
Copper		& $5.8\times 10^7$   \\ \index{copper} % HCP
Gold		& $4.4\times 10^7$   \\ \index{gold} % HCP
Aluminum	& $3.7\times 10^7$   \\ \index{aluminum} % HCP
Iron		& $1.0\times 10^7$   \\ \index{iron} % Wikipedia page "Iron"
Platinum	& $0.9\times 10^7$   \\ \index{platinum} % HCP
Carbon		& $1.3\times 10^5$   \\ \index{carbon} % Wikipedia page "Carbon"
Silicon		& $4.4\times 10^{-4}$\\ \index{silicon} % Wikipedia page "Silicon"
\end{tabular}

{\bf Water} 
\index{water}
exhibits $\sigma$ ranging from 
about $6~\mu$S/m for highly distilled water (thus, a very poor conductor) to % HCP and Wikipedia.
about $5$~S/m for seawater (thus, a relatively good conductor), varying also with temperature and pressure. % Wikipedia
Tap water is typically in the range of 5--50~mS/m, depending on the level of impurities present. %

{\bf Soil} 
\index{soil}
typically exhibits $\sigma$ in the range $10^{-4}$~S/m for dry soil to about $10^{-1}$~S/m for wet soil, varying also due to chemical composition. % Wikipedia page on soil resistivity
  
{\bf Non-conductors}.  Most other materials that are not well-described as conductors or semiconductors and are dry exhibit $\sigma < 10^{-12}$~S/m.  
Most materials that are considered to be \emph{insulators}, including air and common dielectrics, exhibit $\sigma < 10^{-15}$~S/m, often by several orders of magnitude. 

%--------------------------------

{\bf Additional Reading:}
\begin{itemize}
\item \href{http://hbcponline.com}{\emph{CRC Handbook of Chemistry and Physics}}.
\item ``\href{https://en.wikipedia.org/wiki/Conductivity_(electrolytic)}{Conductivity (electrolytic)}'' on Wikipedia.
\item ``\href{https://en.wikipedia.org/wiki/Electrical_resistivity_and_conductivity}{Electrical resistivity and conductivity}'' on Wikipedia.
\item ``\href{https://en.wikipedia.org/wiki/Soil_resistivity}{Soil resistivity}'' on Wikipedia.
\end{itemize}

\index{conductivity!of common materials|)}


